\subsubsection*{Solution}
Quantum mechanics has no genuine single-particle interference beyond second order, although multislit configurations can still witness nonclassicality through suitable probability inequalities; see \href{\detokenize{https://arxiv.org/abs/1403.4147}}{Barnum, Müller, and Ududec} and \href{\detokenize{https://arxiv.org/abs/2003.12114}}{Horvat and Dakić}.

A necessary clarification is that the measurement $\{\Pi_0,\Pi_1\}$ is not specified in the question. The results below use the standard interpretation: for every encoding strategy, one chooses the fixed binary measurement that maximizes the entire inequality. Without this optimization, the answers are measurement-dependent and are not uniquely determined.

\paragraph*{1. Reformulating the violation as an eigenvalue problem}

Let

\[
\rho_0=|\Psi_N\rangle\langle\Psi_N|, \qquad \rho_i=|\psi_i\rangle\langle\psi_i|,
\]

where $\rho_i$ is the state when only the $i$th input equals $1$.

Since the measurement is binary,

\[
\Pi_1=I-\Pi_0.
\]

Therefore,

\[
\adjustbox{max width=\linewidth}{$\displaystyle \begin{aligned}
&\operatorname{Tr}(\Pi_0\rho_0)
+\sum_{i=1}^N\operatorname{Tr}(\Pi_1\rho_i)-N\\
=&\operatorname{Tr}(\Pi_0\rho_0)
+\sum_{i=1}^N\operatorname{Tr}\bigl[(I-\Pi_0)\rho_i\bigr]-N\\
=&\operatorname{Tr}\left[
\Pi_0\left(\rho_0-\sum_{i=1}^N\rho_i\right)
\right].
\end{aligned}$}
\]

Define the Hermitian witness operator

\[
W=\rho_0-\sum_{i=1}^N\rho_i.
\]

The variational form of the Holevo-Helstrom theorem gives

\[
\max_{0\leq \Pi_0\leq I}\operatorname{Tr}(\Pi_0W) =\operatorname{Tr}(W_+),
\]

where $W_+$ is the positive part of $W$. This variational characterization is reviewed, for example, in \href{\detokenize{https://journals.aps.org/prxquantum/pdf/10.1103/PRXQuantum.4.040330}}{PRX Quantum 4, 040330}.

Because $W$ is a rank-one positive operator minus a positive-semidefinite operator, it can have at most one positive eigenvalue. Hence

\[
{\delta=\max\{0,\lambda_{\max}(W)\}}.
\]

The optimal measurement is

\[
\Pi_0=|w_+\rangle\langle w_+|, \qquad \Pi_1=I-\Pi_0,
\]

where $|w_+\rangle$ is the positive-eigenvalue eigenvector of $W$.

\medskip\hrule\medskip

\paragraph*{2. Witness matrix for the phase strategy}

Write $\mathbf 1=(1,\ldots,1)^T$. If only path $i$ receives the phase $\theta_i$, then

\[
|\psi_i\rangle =\frac{1}{\sqrt N} \left[\mathbf 1+\left(e^{i\theta_i}-1\right)|e_i\rangle\right].
\]

Define

\[
d_i=e^{i\theta_i}-1, \qquad d=(d_1,\ldots,d_N)^T, \qquad Q=\operatorname{diag}(|d_1|^2,\ldots,|d_N|^2).
\]

A direct expansion gives

\[
N\sum_{i=1}^N\rho_i =N\mathbf 1\mathbf 1^\dagger +d\mathbf 1^\dagger+\mathbf 1d^\dagger+Q.
\]

Since $N\rho_0=\mathbf 1\mathbf 1^\dagger$, the rescaled witness

\[
K:=NW
\]

is

\[
K=-(N-1)\mathbf 1\mathbf 1^\dagger -d\mathbf 1^\dagger-\mathbf 1d^\dagger-Q.
\]

For the specified phases,

\[
\theta_i=
\begin{cases}
\phi,&i\leq k,\\
\pi,&i=k+1,\\
-\phi,&i>k+1.
\end{cases}
\]

All vectors whose components sum to zero inside either group of $k$ equally phased paths are eigenvectors of $K$ with eigenvalue

\[
-2(1-\cos\phi)\leq 0.
\]

Thus a positive eigenvalue can only occur in a three-dimensional subspace.

Introduce

\[
|L\rangle=\frac1{\sqrt k}\sum_{i=1}^k|e_i\rangle, \qquad |C\rangle=|e_{k+1}\rangle, \qquad |R\rangle=\frac1{\sqrt k}\sum_{i=k+2}^{2k+1}|e_i\rangle,
\]

and the orthonormal basis

\[
|S\rangle=\frac{|L\rangle+|R\rangle}{\sqrt2}, \qquad |A\rangle=\frac{i(|L\rangle-|R\rangle)}{\sqrt2}.
\]

Let

\[
c=\cos\phi,\qquad s=\sin\phi.
\]

In the ordered basis $\{|S\rangle,|C\rangle,|A\rangle\}$,

\[
K_{\mathrm{red}}=
\begin{pmatrix}
X &-\sqrt{2k}(2k+c-3)&-2ks\\
-\sqrt{2k}(2k+c-3)&-2k&-\sqrt{2k}s\\
-2ks&-\sqrt{2k}s&-2(1-c)
\end{pmatrix},
\]

where

\[
X=-4k^2-4kc+4k-2+2c.
\]

The optimized violation is therefore

\[
\delta(\phi)=\frac{1}{2k+1} \max\left\{0,\lambda_{\max}(K_{\mathrm{red}})\right\}.
\]

\medskip\hrule\medskip

\paragraph*{3. Case $k=1$, or $N=3$}

For $k=1$, the characteristic polynomial factorizes as

\[
\det(K_{\mathrm{red}}-\mu I) =(-4-\mu) \left[\mu^2+2\mu-4(1-\cos\phi)\right].
\]

Consequently, the three eigenvalues of $K$ are

\[
\mu_1=-4,
\]

and

\[
\mu_{\pm} =-1\pm\sqrt{5-4\cos\phi}.
\]

For $0\leq\phi\leq\pi$,

\[
\mu_+=-1+\sqrt{5-4\cos\phi}\geq0,
\]

with equality only at $\phi=0$. Since $W=K/3$,

\[
{ \delta(\phi) =\frac{\sqrt{5-4\cos\phi}-1}{3} }.
\]

For example,

\[
\delta(0)=0, \qquad \delta(\pi)=\frac{\sqrt9-1}{3}=\frac23.
\]

\medskip\hrule\medskip

\paragraph*{4. Range of phases giving a violation}

A violation begins when the largest eigenvalue passes through zero. Thus set $\mu=0$ in the reduced characteristic polynomial. Direct evaluation gives

\[
\det K_{\mathrm{red}} =-16k(1-c)\left[(k-2)-(k-1)c\right].
\]

Because $K$ can have at most one positive eigenvalue, away from zero eigenvalues a positive determinant of the three-dimensional block is equivalent to the existence of one positive eigenvalue.

\subparagraph*{Case $k=1$}

The determinant becomes

\[
\det K_{\mathrm{red}}=16(1-c).
\]

It is positive for every $c<1$, or equivalently every $\phi>0$. Therefore,

\[
{T_1=(0,\pi]}.
\]

\subparagraph*{Case $k\geq2$}

Since $1-c>0$ for $\phi>0$, positivity requires

\[
(k-2)-(k-1)c<0.
\]

Thus

\[
c>\frac{k-2}{k-1}.
\]

Because $\cos\phi$ is strictly decreasing on $[0,\pi]$,

\[
0<\phi< \arccos\left(\frac{k-2}{k-1}\right).
\]

The boundary is excluded because the positive eigenvalue vanishes there. Hence

\[
{ T_k= \left( 0,\, \arccos\left(\frac{k-2}{k-1}\right) \right), \qquad k\geq2. }
\]

For example,

\[
T_2=(0,\pi/2), \qquad T_3=(0,\pi/3).
\]

For $\det K_{\mathrm{red}}=0$, one can immediately verify that no violation occurs.

\medskip\hrule\medskip

\paragraph*{5. Phase giving the maximal violation}

Let $\mu(c)$ denote the positive eigenvalue of $K_{\mathrm{red}}$. Define

\[
F_k(\mu,c)=\det(K_{\mathrm{red}}-\mu I).
\]

For reference, with

\[
\begin{aligned}
x&=-4k^2+4k-2-\mu-(4k-2)c,\\
y&=-2k-\mu,\\
z&=-2(1-c)-\mu,\\
h&=2k+c-3,
\end{aligned}
\]

the determinant can be written as

\[
\begin{aligned}
F_k(\mu,c)
={}&xyz
-8k^2h(1-c^2)
-2kx(1-c^2)\\
&-4k^2y(1-c^2)
-2kzh^2.
\end{aligned}
\]

Although individual terms appear cubic in $c$, the cubic terms cancel, so $F_k$ is quadratic in $c$.

At an interior extremum,

\[
F_k(\mu,c)=0.
\]

Implicit differentiation gives

\[
\frac{d\mu}{dc} =-\frac{\partial F_k/\partial c} {\partial F_k/\partial\mu}.
\]

Thus an interior maximum obeys

\[
F_k(\mu,c)=0, \qquad \frac{\partial F_k}{\partial c}=0.
\]

For $k\geq2$, solving these two algebraic equations gives the unique solution inside the violating interval:

\[
\mu_{\max} =\frac{1}{2(k-1)(2k-1)}
\]

and

\[
\cos\phi_{\max} = 1- \frac{ 16k^3-28k^2+14k-1 }{ 4(k-1)^2(2k-1)(4k-1) }.
\]

Equivalently,

\[
\cos\phi_{\max} = \frac{ 32k^4-104k^3+112k^2-46k+5 }{ 4(k-1)^2(2k-1)(4k-1) }.
\]

Therefore,

\[
{ \phi_{\max} = \arccos\left[ 1- \frac{ 16k^3-28k^2+14k-1 }{ 4(k-1)^2(2k-1)(4k-1) } \right], \qquad k\geq2. }
\]

The corresponding maximum violation is a useful consistency check:

\[
\delta_{\max} =\frac{\mu_{\max}}{N} = { \frac{1} {2(k-1)(2k-1)(2k+1)} }, \qquad k\geq2.
\]

For $k=1$, $\delta(\phi)$ increases monotonically with $\phi$, because

\[
\frac{d\delta}{d\phi} = \frac{2\sin\phi} {3\sqrt{5-4\cos\phi}} \geq0.
\]

Hence

\[
{\phi_{\max}=\pi,\qquad \delta_{\max}=\frac23} \qquad (k=1).
\]

One can verify that these stationary points are global maxima.
